\chapter{Structurer, Sales and Sales-Trader}\label{ch:in:structurer-sales-and-sales-trader}

The European trade association of structured-product issuers publishes each quarter what its members' markets hold and
trade. At the end of March 2026 investors in six European markets held 487 billion euros of structured investment
products issued as securities, most of them in Switzerland and Germany, and 163\,413 new investment products were
listed in the first quarter alone. Behind each note is a \omterm{def:in:structurer-sales-and-sales-trader:structurer}{structurer} who designed it, a salesperson who placed it, a
trading desk that hedges it, and a margin that the note's buyer pays and the bank's desk earns. This chapter describes
the three client-facing roles in which quantitative skill is part of the job, computes the margin in a typical note,
and sets out the rules that bind people who sell financial products.

\begin{center}\footnotesize
\begin{tabular}{@{}p{2.4cm}p{3.3cm}p{3.3cm}p{3.3cm}@{}}
\toprule
\multicolumn{4}{@{}l}{\textbf{Role cards: the client-facing roles}}\\
\midrule
 & \omterm{def:in:structurer-sales-and-sales-trader:structurer}{structurer} & \omterm{def:in:structurer-sales-and-sales-trader:sales}{institutional salesperson} & sales-trader\\
\midrule
works on & product design and pricing & client relationships and ideas & client orders and their execution\\
horizon & weeks to years (the note's life) & months to years & seconds to days\\
credited with & the desk's margin & client revenue & client flow\\
codes in filings & 13-2099.01, 13-2051, 41-3031 & 41-3031 & 41-3031, 13-2099.01\\
taught in & Book~5, ch.~18--19 & Book~9, ch.~24 & Book~1, ch.~2; Book~10, ch.~20\\
\bottomrule
\end{tabular}
\end{center}

\section{The structurer}

\begin{definition}[Structurer]\label{def:in:structurer-sales-and-sales-trader:structurer}
A \emph{structurer}\index{structurer} is a quantitative specialist on a bank's markets side who designs structured products for clients
and distributors: chooses the payoff, the underlying, the protection and the coupon, prices them with the desk's models,
and sets the structuring margin within the product's approval and the trading desk's capacity to hedge it.
\end{definition}

A \omterm{def:in:structurer-sales-and-sales-trader:structurer}{structurer} works between three parties. The client, or the private bank or retail network that sells to clients, wants
a coupon or a protection level; the trading desk that will hedge the note wants risks it can hold (Book~5, chapter~27;
Book~9, chapter~28); the bank wants a margin. Book~5, chapter~19, describes the business, and its chapter~18 the
autocallable, the payoff this chapter prices. The \omterm{def:in:structurer-sales-and-sales-trader:structurer}{structurer}'s quantitative work is pricing and
its sensitivities, understanding which parameters the price depends on and which the client does not see.

\begin{definition}[Product governance]\label{def:in:structurer-sales-and-sales-trader:governance}
\emph{Product governance}\index{product governance} is the set of processes by which a firm that manufactures or
distributes financial instruments approves each product for an identified target market of end clients, checks that the
distribution strategy suits that market, and reviews the product during its life.
\end{definition}

The \omterm{def:in:structurer-sales-and-sales-trader:structurer}{structurer}'s product must pass the firm's product approval process before it is sold: the European rules require a
manufacturer to ``maintain, operate and review a process for the approval of each financial instrument'' that specifies
``an identified target market of end clients''. The target market is a statement about who should buy the note; the
\omterm{def:in:structurer-sales-and-sales-trader:structurer}{structurer} writes much of it.

\section{The margin in the note}

\begin{method}[Structuring margin of a note]\label{met:in:structurer-sales-and-sales-trader:margin}
A note sold at par, 100, has fair value $V(c)$ per 100 under the desk's model for coupon $c$; its margin is $100-V(c)$.
For a Phoenix autocallable the note's cash flows are linear in the coupon, so along the same simulated paths
$V(c)=V(0)+c\,\partial V/\partial c$, and the coupon that leaves a target margin $m$ is
$c_m=\bigl(100-m-V(0)\bigr)/(\partial V/\partial c)$.
\end{method}

\begin{example}[A five-year autocallable]\label{ex:in:structurer-sales-and-sales-trader:note}
An illustrative note on one equity index: five years, annual observations, autocall when the index is at or above its
initial level, a coupon of 6 per 100 a year paid when the index is at or above 70\% (with memory), capital repaid unless
the index ends below 60\% of its initial level, when the holder takes the loss. Dividend yield 2\%, flat volatility,
200\,000 paths. At 6\% the margin is $-1.65$ per 100 with rates at 1\% and volatility 20\% (the bank would lose), 1.83
with rates at 4\% and volatility 20\%, 6.91 at 1\% and 30\%, and 9.13 at 4\% and 30\%. The coupon that leaves a 2\%
margin is 4.23\% at 1\% and 20\%, 5.91\% at 4\% and 20\%, 8.79\% at 1\% and 30\%, and 10.22\% at 4\% and 30\%
(\cref{fig:in:structurer-sales-and-sales-trader:coupon}).
\end{example}

\begin{center}\footnotesize
\begin{tabular}{@{}lrrrr@{}}
\toprule
 rates; volatility & 1\%; 20\% & 4\%; 20\% & 1\%; 30\% & 4\%; 30\%\\
\midrule
margin at a 6\% coupon, per 100 & $-1.65$ & 1.83 & 6.91 & 9.13\\
coupon for a 2\% margin, \% a year & 4.23 & 5.91 & 8.79 & 10.22\\
coupon for no margin, \% a year & 5.20 & 6.97 & 9.92 & 11.41\\
autocalled at year one, \% & 44.0 & 50.0 & 42.7 & 46.7\\
capital loss at maturity, \% & 15.8 & 9.4 & 25.1 & 19.3\\
\bottomrule
\end{tabular}
\end{center}

Two lessons are the \omterm{def:in:structurer-sales-and-sales-trader:structurer}{structurer}'s daily material. Higher rates buy a higher coupon for the same margin: the note is a
bond plus options sold by the investor, and a bond discounted at 4\% is cheaper than at 1\%. Higher volatility buys a
higher coupon too, because the investor is short the downside and paid for it; but the chance of losing capital rises
with it, from 15.8\% to 25.1\% of paths at 1\% rates. A client comparing coupons across notes compares, without seeing
it, the volatility sold.

\begin{omfigure}
\begin{tikzpicture}
\begin{axis}[ybar,area legend,width=10.5cm,height=5.4cm,xmin=0.5,xmax=4.5,ymin=0,ymax=13,bar width=14pt,xtick={1,2,3,4},
  xticklabels={{1\%, 20\%},{4\%, 20\%},{1\%, 30\%},{4\%, 30\%}},xlabel={\small rates, volatility},
  ylabel={\small coupon, \% a year},tick label style={font=\footnotesize},grid=major,grid style={gray!20},
  table/col sep=comma,legend style={font=\scriptsize,at={(0.03,0.97)},anchor=north west,legend cell align=left}]
\addplot[fill=omDef,draw=none] table[x=pos,y=fair]{figdata/industry/23-structurer-sales-and-sales-trader/coupon.csv};
\addplot[fill=omDef!35,draw=none] table[x=pos,y=zero]{figdata/industry/23-structurer-sales-and-sales-trader/coupon.csv};
\legend{coupon leaving a 2\% margin,coupon leaving no margin}
\end{axis}
\end{tikzpicture}

\omcaption{The coupon of the chapter's illustrative five-year autocallable that leaves a 2\% margin and no margin, at two
rate levels and two volatilities. The gap between each pair is the margin expressed as coupon. Data:
\texttt{in\_structurer.margins}, from \texttt{firm.roles.structuring\_margin} on \texttt{firm.autocall}.}\label{fig:in:structurer-sales-and-sales-trader:coupon}
\end{omfigure}

\section{Institutional sales and coverage}

\begin{definition}[Institutional salesperson]\label{def:in:structurer-sales-and-sales-trader:sales}
An \emph{institutional salesperson}\index{institutional salesperson} is a person who covers a set of professional
clients (asset managers, hedge funds, pension funds, insurers, corporates) for a bank or broker: brings them the firm's
ideas, prices and products, takes their interest to the desks, and is credited with the revenue the clients generate.
\end{definition}

The salesperson owns the relationship and the franchise that client tiering of Book~9, chapter~24, formalises. Where
markets trade electronically, the quantitative side of the job is the client's execution analysis, the desk's pricing and
the risk a trade leaves with the bank, and a salesperson who can discuss all three is more useful to the client. The role is
judged on client revenue and share of the client's business, which the bank measures and the salesperson's pay follows
(chapter~13).

\section{The sales-trader}

The sales-trader of Book~1, chapter~2, stands between clients' orders and the market: takes a \omterm{def:in:portfolio-manager-and-pod-analyst:roles}{portfolio manager}'s order,
decides with the client how to work it, runs it through the firm's algorithms or its traders, and reports the result
against the client's benchmark. The high-touch trading of Book~10, chapter~20, is this job. Its quantitative content is
execution: transaction cost analysis, the choice of algorithm and venue, the market impact of large orders (Book~10), and
the sales-trader who can explain a fill against arrival price keeps the client.

\section{Quantitative skills in client-facing roles}

The three roles use quantitative skill differently.
\begin{itemize}
\item \textbf{The \omterm{def:in:structurer-sales-and-sales-trader:structurer}{structurer}} prices and designs: stochastic calculus and Monte Carlo (Book~5), the sensitivities of a
payoff to volatility, correlation and rates, and the hedging cost the desk will bear.
\item \textbf{The salesperson} translates: turns a desk's view or a product into a client's problem and back, and needs
enough of the model to know when a price is wrong.
\item \textbf{The sales-trader} measures: execution quality against benchmarks, in numbers the client can check.
\end{itemize}
In all three, the audience is a client who may know less, and the skill includes saying what the client is buying.

\section{The rules for client-facing work}

People who sell financial products are licensed and supervised in ways that quants inside a trading firm are not.

\begin{dated}{2026-09}{Rules for people who design and sell products}\label{dat:in:structurer-sales-and-sales-trader:rules}
European Union: MiFID II (Directive 2014/65/EU) Article 16(3) requires a manufacturer's product approval process with ``an
identified target market of end clients''; Article 24(2) requires products ``designed to meet the needs of an identified
target market'' and a compatible distribution strategy; Article 25(1) requires staff who advise or inform clients to
``possess the necessary knowledge and competence''. ESMA's guidelines on MiFID II \omterm{def:in:structurer-sales-and-sales-trader:governance}{product governance} (ESMA35-43-3448, 3
August 2023) set out how Articles 16(3) and 24(2) apply. United States: a general securities representative passes
FINRA's Series 7, which ``assesses the competency of an entry-level registered representative'', with the Securities
Industry Essentials exam as a corequisite; the Series 63, a NASAA exam administered by FINRA, covers state securities law.
\end{dated}

For a quant moving into a client-facing role, the rules change the job in three ways. The product must be approved
before it is sold and reviewed while it lives; the person must be qualified for the activity; and the firm's records
include ``the recording of telephone conversations or electronic communications'' relating to client orders (MiFID II,
Article 16(7)). Chapter~24 covers the compliance function that enforces them.

\begin{dated}{2026-03}{The market for structured products}\label{dat:in:structurer-sales-and-sales-trader:market}
EUSIPA's report for the first quarter of 2026: sales of investment and leverage products on trading venues in the
reporting markets, 74 billion euros (investment products 19 billion); outstanding volume of structured investment products
at the end of March, 486\,635 million euros (Switzerland 286\,264 million, Germany 99\,890, Italy 66\,788, Austria
17\,462, Belgium 11\,678, Luxembourg 4\,553); 465\,721 investment products listed; 163\,413 launched in the quarter.
\end{dated}

\begin{omfigure}
\begin{tikzpicture}
\begin{axis}[xbar,width=10.5cm,height=4.8cm,xmin=0,xmax=330,ymin=0.4,ymax=6.6,bar width=9pt,ytick=data,
  yticklabels from table={figdata/industry/23-structurer-sales-and-sales-trader/eusipa.csv}{market},
  yticklabel style={font=\footnotesize},xticklabel style={font=\footnotesize},
  xlabel={\small outstanding structured investment products, end-March 2026, EUR billion},grid=major,grid style={gray!20},
  table/col sep=comma,nodes near coords={\pgfplotspointmeta\%},nodes near coords style={font=\scriptsize,anchor=west,
  xshift=1pt},
  point meta=explicit symbolic]
\addplot[fill=omDef,draw=none] table[x=bn,y=pos,meta=share]{figdata/industry/23-structurer-sales-and-sales-trader/eusipa.csv};
\end{axis}
\end{tikzpicture}

\omcaption{Outstanding volume of structured investment products issued as securities in six European markets at the end
of March 2026, EUSIPA's report; labels give each market's share of the total. Data: EUSIPA Market Report Q1 2026,
through \texttt{in\_structurer.eusipa\_shares}.}\label{fig:in:structurer-sales-and-sales-trader:eusipa}
\end{omfigure}

The market is concentrated: Switzerland holds 58.8\% of the six markets' outstanding volume and Germany 20.5\%
(\cref{fig:in:structurer-sales-and-sales-trader:eusipa}). A desk that issues one to five billion euros of notes a year
at a 2\% margin earns 20 to 100 million euros of structuring revenue; spread over ten \omterm{def:in:structurer-sales-and-sales-trader:structurer}{structurers}, 2 to 10 million euros
each. Both the issuance and the team are illustrative; the arithmetic shows why a \omterm{def:in:structurer-sales-and-sales-trader:structurer}{structurer}'s pay follows the margin
and the volume the desk can place.

\begin{dated}{2025-05}{What sales agents are paid: the survey}\label{dat:in:structurer-sales-and-sales-trader:pay}
Occupational survey, May 2025, securities, commodities and financial services sales agents (41-3031), which includes
institutional salespeople, sales-traders and brokers: securities industry 173\,040 employed, median \$103\,030, 10th--90th
percentile \$55\,130--309\,440; banks 257\,110 employed, median \$61\,440. The filings' \omterm{def:in:structurer-sales-and-sales-trader:structurer}{structurer} and sales titles are too
few to publish separately. Wages only.
\end{dated}

\begin{omfigure}
\begin{tikzpicture}
\begin{axis}[width=10.5cm,height=4.4cm,xmin=0,xmax=330,ymin=0.4,ymax=3.6,ytick=data,
  yticklabels from table={figdata/industry/23-structurer-sales-and-sales-trader/survey.csv}{label},
  yticklabel style={font=\footnotesize},xticklabel style={font=\footnotesize},xlabel={\small annual wage, \$ thousand},
  grid=major,grid style={gray!20},table/col sep=comma]
\addplot[draw=none,mark=none,error bars/.cd,x dir=both,x explicit,error mark=none,error bar style={omDef,line width=1pt}]
  table[x=p50,y=pos,x error plus expr=\thisrow{p90}-\thisrow{p50},x error minus expr=\thisrow{p50}-\thisrow{p10}]
  {figdata/industry/23-structurer-sales-and-sales-trader/survey.csv};
\addplot[draw=none,mark=none,error bars/.cd,x dir=both,x explicit,error mark=none,error bar style={omDef!40,line width=7pt}]
  table[x=p50,y=pos,x error plus expr=\thisrow{p75}-\thisrow{p50},x error minus expr=\thisrow{p50}-\thisrow{p25}]
  {figdata/industry/23-structurer-sales-and-sales-trader/survey.csv};
\addplot[only marks,mark=|,mark size=5pt,line width=1.5pt,omThm] table[x=p50,y=pos]{figdata/industry/23-structurer-sales-and-sales-trader/survey.csv};
\end{axis}
\end{tikzpicture}

\omcaption{Wages of securities, commodities and financial services sales agents (41-3031) by industry, May 2025
occupational survey: 10th to 90th percentile (thin), interquartile range (thick), median (mark). The occupation includes
retail brokers as well as institutional sales. Data: \texttt{data/industry/oews\_roles.csv}, through
\texttt{in\_structurer.survey}.}\label{fig:in:structurer-sales-and-sales-trader:survey}
\end{omfigure}

\section{Tutorial: the margin in the note}\label{tut:in:structurer-sales-and-sales-trader:tut}

\begin{tutorial}
\textbf{Goal.} Price an autocallable, find its margin and the coupon that leaves a target margin, and put the desk's
revenue per \omterm{def:in:structurer-sales-and-sales-trader:structurer}{structurer} beside the market's size.
\textbf{End state:} \cref{fig:in:structurer-sales-and-sales-trader:coupon,fig:in:structurer-sales-and-sales-trader:eusipa}
and the margin table.
\begin{tutsteps}
\item \textbf{Term sheet.} Book~5's \texttt{firm.autocall.TermSheet} with annual observations, a Phoenix coupon with
memory and protection at maturity.
\item \textbf{Value and margin.} \texttt{firm.roles.note\_value} simulates the index and discounts the cash flows;
\texttt{structuring\_margin} uses the note's linearity in the coupon to find the coupon for any margin
(\cref{lst:in:structurer-sales-and-sales-trader:margin}).
\omcode{code/firm/roles/firm_roles.py}{408}{416}{The structuring margin and the coupon that leaves a target margin.}\label{lst:in:structurer-sales-and-sales-trader:margin}
\item \textbf{Grid.} Two rate levels and two volatilities, the same paths for each coupon.
\item \textbf{Scale.} Revenue per \omterm{def:in:structurer-sales-and-sales-trader:structurer}{structurer} from an illustrative issuance range and team.
\end{tutsteps}
\end{tutorial}

For the example's grid the table above results; at the 6\% coupon, moving rates from 1\% to 4\% adds 3.48 points of
margin at 20\% volatility.

\textbf{What to change next.} Add the issuer's funding spread, which the note's buyer also pays; price under local
volatility (Book~5, chapter~18) and compare; make the note a worst-of on three indices and watch the coupon and the loss
probability rise together.

\section{Build: client-facing cards and the structuring margin}\label{bld:in:structurer-sales-and-sales-trader:roles}

\begin{build}
\bfield{Purpose} Add the client-facing roles to the registry and compute the margin in a note.
\bfield{Interface} \texttt{firm.roles}: the cards \texttt{structurer}, \texttt{institutional salesperson},
\texttt{sales-trader}; \texttt{note\_value(sheet, r, q, vol, n, seed, steps)};
\texttt{structuring\_margin(make\_sheet, coupon, target, r, q, vol, n, seed, steps)}; on \texttt{firm.autocall}.
\bfield{Rules} Values per 100 sold at par; the same random numbers for every coupon, so that the coupon for a margin is
exact along the paths; volatility flat unless the caller passes a function.
\bfield{Acceptance tests} \texttt{code/firm/roles/tests/}: the note repriced at the returned coupon has exactly the target
margin; the coupon for no margin is above the coupon for a positive one.
\bfield{Stretch} Worst-of notes; the margin's sensitivity to each parameter; a term sheet generator that searches the
protection level for a coupon target.
\end{build}

\begin{omsources}
\item EUSIPA, Market Report on Structured Investment and Leverage Products Q1/2026.
\item Directive 2014/65/EU (MiFID II), Articles 16, 24 and 25; ESMA35-43-3448, Guidelines on MiFID II product governance.
\item FINRA, Series 7 and Series 63 qualification exams.
\item Bureau of Labor Statistics, occupational survey, May 2025; Book~5, chapters~18 and~19.
\end{omsources}

\section{Exercises}

\begin{exercise}[$\star$]\label{exo:in:structurer-sales-and-sales-trader:1}
What share of the six markets' outstanding structured investment products did Switzerland and Germany hold together at
the end of March 2026?
\end{exercise}

\begin{exercise}[$\star$]\label{exo:in:structurer-sales-and-sales-trader:2}
At rates of 4\% and volatility of 20\%, what is the margin of the 6\% note, and what coupon would leave no margin?
\end{exercise}

\begin{exercise}[$\star$]\label{exo:in:structurer-sales-and-sales-trader:3}
Which article of MiFID II requires a manufacturer's product approval process, and what must it specify?
\end{exercise}

\begin{exercise}[$\star\star$]\label{exo:in:structurer-sales-and-sales-trader:4}
Why does a higher interest rate allow a higher coupon for the same margin?
\end{exercise}

\begin{exercise}[$\star\star$]\label{exo:in:structurer-sales-and-sales-trader:5}
A client prefers the 10.22\% note (rates 4\%, volatility 30\%) to the 5.91\% note (rates 4\%, volatility 20\%). What is
he buying with the extra coupon?
\end{exercise}

\begin{exercise}[$\star\star$]\label{exo:in:structurer-sales-and-sales-trader:6}
Why is a \omterm{def:in:structurer-sales-and-sales-trader:structurer}{structurer}'s pay tied to the desk's margin and volume rather than to the note's later performance?
\end{exercise}

\begin{exercise}[$\star\star\star$]\label{exo:in:structurer-sales-and-sales-trader:7}
\emph{Coding.} Price the 6\% note at rates of 4\% and volatility of 20\% with 50\,000 paths and seeds 1 to 20. What is the
standard deviation of the margin across seeds, and is the margin's sign in any doubt?
\end{exercise}

\begin{exercise}[$\star\star\star$]\label{exo:in:structurer-sales-and-sales-trader:8}
\emph{Find the flaw.} ``The note pays 10\% a year and protects capital down to 60\%, so it is safer than holding the
index and pays more than its dividends.''
\end{exercise}

\section{Problem: The Margin in the Note}

\begin{problem}[{Weekend problem --- the margin in the note}]\label{pb:in:structurer-sales-and-sales-trader:1}
A quantitative analyst is offered a move to a structuring desk and wants to understand where its revenue comes from and
what the job's rules are.

\medskip\noindent\textbf{Part I --- The roles.}
\begin{enumerate}
\item Define the \omterm{def:in:structurer-sales-and-sales-trader:structurer}{structurer}, the \omterm{def:in:structurer-sales-and-sales-trader:sales}{institutional salesperson} and \omterm{def:in:structurer-sales-and-sales-trader:governance}{product governance}.
\item Whom does a \omterm{def:in:structurer-sales-and-sales-trader:structurer}{structurer} work between?
\item What does a sales-trader do, and what is high-touch trading?
\item What quantitative skill does each role use?
\item Which book teaches the autocallable?
\end{enumerate}

\medskip\noindent\textbf{Part II --- The margin.}
\begin{enumerate}[resume]
\item State the note of the example.
\item Define the structuring margin and derive the coupon for a target margin.
\item Give the margin at 6\% in the four cases.
\item Give the coupon for a 2\% margin in the four cases.
\item Why does volatility raise the coupon, and what else does it raise?
\end{enumerate}

\medskip\noindent\textbf{Part III --- The market and the rules.}
\begin{enumerate}[resume]
\item Give the market's size and its concentration by country.
\item What revenue per \omterm{def:in:structurer-sales-and-sales-trader:structurer}{structurer} does the example's desk earn?
\item State the MiFID II \omterm{def:in:structurer-sales-and-sales-trader:governance}{product governance} and competence rules.
\item What do FINRA's Series 7 and 63 cover?
\item What do the survey's wages show?
\end{enumerate}

\medskip\noindent\textbf{Part IV --- The verdict.}
\begin{enumerate}[resume]
\item State the \emph{named result}: the structuring margin of the chapter's autocallable at 1\% and 4\% rates, and the
revenue per \omterm{def:in:structurer-sales-and-sales-trader:structurer}{structurer} for a desk's issuance in the chapter's range.
\item What would a funding spread change?
\item What does a \omterm{def:in:structurer-sales-and-sales-trader:structurer}{structurer} owe the end client?
\item What would the analyst's quantitative skills add to the desk, and what would they not?
\item In two sentences, advise the analyst.
\end{enumerate}
\end{problem}

\section{Interview questions}

\begin{interviewq}[$\star$ \iqroles{bank}]\label{iq:in:structurer-sales-and-sales-trader:1}
Decompose a capital-protected note into a bond and an option.
\end{interviewq}

\begin{interviewq}[$\star$ \iqroles{bank, trader}]\label{iq:in:structurer-sales-and-sales-trader:2}
A client asks why this year's autocallable pays more than last year's. Give three reasons.
\end{interviewq}

\begin{interviewq}[$\star\star$ \iqroles{bank}]\label{iq:in:structurer-sales-and-sales-trader:3}
Is the investor in an autocallable long or short volatility, and long or short correlation in a worst-of?
\end{interviewq}

\begin{interviewq}[$\star\star$ \iqroles{bank, trader}]\label{iq:in:structurer-sales-and-sales-trader:4}
A client's 5\% participation order finished 20 basis points worse than arrival. How do you explain it?
\end{interviewq}

\begin{interviewq}[$\star\star$ \iqroles{bank, risk}]\label{iq:in:structurer-sales-and-sales-trader:5}
What goes into a target-market statement for a note like the chapter's?
\end{interviewq}

\begin{interviewq}[$\star\star\star$ \iqroles{bank, researcher}]\label{iq:in:structurer-sales-and-sales-trader:6}
The desk has sold many autocallables on the same index. What risk does the desk hold, and when does it hurt?
\end{interviewq}
